% problem-000442 4 铅铋镍合金EDTA分析
\section*{4. 铅铋镍合金EDTA分析}
某铅合⾦中含有 Pb、Bi、Ni 等元素，称取此合⾦试样 2.420 g，⽤ $\mathrm { H N O _ { 3 } }$ 溶解并定容⾄ $2 5 0 ~ \mathrm { m L } _ { \circ }$ 移取50.00 mL上述试液于250 mL锥形瓶中，调节 $\mathrm { p H } = 1$ ，以⼆甲酚橙为指⽰剂，⽤ $0 . 0 7 5 0 0 \mathrm { m o l } \mathrm { L } ^ { - 1 }$ EDTA 标准溶液滴定，消耗 $5 . 2 5 \mathrm { m L } _ { \mathrm { o } }$ 然后，⽤六次甲基四胺缓冲溶液将pH调⾄5，再以上述EDTA标准溶液滴定，消耗28.76 mL。加⼊邻⼆氮菲，置换出镍配合物中的EDTA，⽤ $0 . 0 4 5 0 0 \mathrm { m o l L ^ { - 1 } P b } ( \mathrm { N O _ { 3 } } )$ _{2}标准溶液滴定置换出的EDTA，消耗 $8 . 7 6 \mathrm { m L } _ { \circ }$ 计算此合⾦试样中Pb、Bi、Ni的质量分数。 $( \mathrm { l g } K _ { \mathrm { B i Y } } = 2 7 . 9 4 , \mathrm { l g } K _ { \mathrm { P b Y } } = 1 8 . 0 4 , \mathrm { l g } K _ { \mathrm { N i Y } } =$ 18.62)
